% ==================================================================
% id:       seq-001
% topic:    סדרות
% subtopic: הגדרת הגבול (ε, N)
% solution: full
% ==================================================================
\textbf{בעזרת הגדרת הגבול, הוכיחו כי:}
\begin{enumerate}\bMs{2}
    \item $\ds{\lim_{n\to\infty}\dfrac{2n + 1}{n-2} =2}$
    
\hint{פשטו את המרחק בין איבר הסדרה לגבול לשבר אחד, ובדקו עבור אילו $n$ הוא קטן מ-$\varepsilon$.}
\begin{solbox}
    נקבע
    $\varepsilon>0$
    ונחפש
    $N$
    טבעי כך ש-
    $\left|\dfrac{2n + 1}{n-2}-2\right|<\varepsilon$
    לכל
    $n>N$.
    נפשט
    $$,\left|\dfrac{2n + 1}{n-2}-2\right|=\left|\dfrac{5}{n-2}\right|$$
    נניח כי
    $n>2$,
    אז די להבין מתי מתקיים האי-שיוויון
    $\frac{5}{n-2}<\varepsilon$.
    האי-שיוויון האחרון שקול ל-
    $\frac{n-2}{5}>\frac{1}{\varepsilon}$,
    ולאחר פישוט אלגברי נקבל
    $$.n>2+\frac{5}{\varepsilon}$$
לכן נבחר
    $N=\lfloor 2+\frac{5}{\varepsilon}\rfloor$
    ואז לכל
    $n>N$
    מתקיים
    $\left|\dfrac{2n + 1}{n-2}-2\right|<\varepsilon$.
\end{solbox}

    \item $\ds{\lim_{n\to\infty}\left(
    17+\frac{(-1)^n}{n}\right)=17}$

    \begin{solbox}
    נסמן
    $a_n=
    17+\frac{(-1)^n}{n}$.
       יהי
       $\varepsilon >0$.
       נסתכל על הביטוי
       $$|a_n-17|=\left|
       \left(
    17+\frac{(-1)^n}{n}\right)-17
       \right|=\left|
       \frac{(-1)^n}{n}
       \right|=\frac{1}{n}$$
       נרצה כי זה יהיה קטן מ-$\varepsilon$,
       כלומר נרצה שיתקיים
       $\frac{1}{n}<\varepsilon$,
       שזה שקול לכך שמתקיים
       $n>\frac{1}{\varepsilon}$.
       לכן נבחר
       $N=\lfloor\frac{1}{\varepsilon}\rfloor +1$
       ואז לכל
       $n>N$
       מתקיים
       $|a_n-17|<\varepsilon$,
       כמבוקש.
\end{solbox}

    \item
    $\ds{\lim_{n\to\infty}\dfrac{n + \sin(n)}{n^2+3n} = 0}$
      
\hint{חסמו את הביטוי מלמעלה בביטוי פשוט יותר: השתמשו ב-$|\sin n|\le1$ והקטינו את המכנה.}
\begin{solbox}
        נסמן
        $$.a_n=\dfrac{n + \sin(n)}{n^2+3n}$$
    נקבע
    $\varepsilon>0$
    ונחפש
    $N$
    כך ש-
    $\left|a_n-0\right|<\varepsilon$
    לכל
    $n>N$,
    כלומר,
    $$.\left|\dfrac{n + \sin(n)}{n^2+3n}\right|<\varepsilon$$
    זהו אי-שיוויון מאתגר, ננסה "להקל על חיינו", נשים לב כי
    $$,\left|\dfrac{n + \sin(n)}{n^2+3n}\right|\leq\dfrac{|n| + |\sin(n)|}{n^2+3n}\leq\dfrac{n + 1}{n^2+3n}\leq\dfrac{n + 1}{n^2}\leq\frac{2n}{n^2}=\frac{2}{n}$$
    כאשר המעבר הראשון מבוסס עך אי-שיוויון המשולש ואילו המעבר השני על העובדה כי
    $|\sin x|\leq 1$
    לכל
    $x\in\mathbb{R}$.
    לכן נחפש
    $N$
עבורו
    $\frac{2}{n}<\varepsilon$
    לכל
    $n>N$.
    נבחר
    $N=\lfloor\frac{2}{\varepsilon}\rfloor$,
    אזי לכל
    $n>N$
    $\left|a_n-0\right|<\varepsilon$,
    מה שמוכיח את הנטען.
\end{solbox}

    \item
$\displaystyle\lim_{n \to \infty}\frac{5n-12}{2n+3}=\frac{5}{2}$.

\hint{פשטו את המרחק בין איבר הסדרה ל-$\frac52$ לשבר אחד. כדי לקבל בחירה פשוטה של $N$ אפשר לחסום אותו מלמעלה בביטוי מהצורה $\frac{c}{n}$.}
\begin{solbox}
    נסמן
    $$a_n=\dfrac{5n-12}{2n+3}$$
    ונקבע
    $\varepsilon>0$.
    נתבונן בהפרש הבא
    $$\left|a_n-\frac{5}{2}\right|=\left|\dfrac{5n-12}{2n+3}-\frac{5}{2}\right|=\left|
    \frac{10n-24-10n-15}{2(2n+3)}
    \right|=\frac{39}{4n+6}$$
    נרצה כי הביטוי האחרון יהיה קטן מ-$\varepsilon$,
    נוכל לפשט זאת כך
    $$\frac{39}{4n+6}<\varepsilon \iff{0} \frac{1}{4}\cdot\left(\frac{39}{\varepsilon}-6\right)<n$$
    קיבלנו דרישה מהסוג שאנו אוהבים, אך יש לשים לב כי המספר ממנו
    $n$
    צריך להיות גדול יותר הפעם יכול להיות שלילי (למשל, אם
    $\varepsilon=7$), ובמקרה כזה הדרישה הזו מתקיימת לכל ה-$n$-ים, לכן נבחר
    $$N=\max\{\lfloor \frac{1}{4}\cdot\left(\frac{39}{\varepsilon}-6\right)\rfloor ,1 \}$$
ואז בוודאות לכל
$n>N$
מתקיים
$|a_n-5/2|<\varepsilon$,
כמבוקש.

שימו לב כי ניתן היה להמשיך בדרך קצת אחרת ולומר כי
$$\frac{39}{4n+6}\le \frac{39}{4n}\le\frac{40}{4n}=\frac{10}{n}$$
דבר המפשט את הבחירה להיות
$N=\lfloor \frac{10}{\varepsilon}\rfloor$,
בדומה לתרגילים הקודמים.
\end{solbox}
    \item $\displaystyle\lim_{n \to \infty}\frac{5n-12}{2n+3}\neq\frac{3}{2}$
    
\hint{נסחו את שלילת הגדרת הגבול, ופשטו את המרחק בין איבר הסדרה ל-$\frac32$: האם הוא נשאר רחוק מאפס עבור $n$ גדולים?}
\begin{solbox}
     דרך א׳: נניח בשלילה כי הגבול כן שווה
     $\frac{3}{2}$,
     לכן
לכל
$\varepsilon>0$
קיים
$N\in\mathbb{N}$
כך שלכל
$n>N$
מתקיים
$$.\left|\frac{5n-12}{2n+3}-\frac{3}{2}\right|<\varepsilon$$
נפשט את את האי-שיוויון ונקבל
\begin{align*}
\left|\frac{5n-12}{2n+3}-\frac{3}{2}\right|<\varepsilon
&\iff{1}
\left|\frac{4n-33}{4n+6}\right|<\varepsilon
\\&\iff{1}
-\varepsilon<\frac{4n-33}{4n+6}<\varepsilon
\\&\iff{1}
-4\varepsilon n-6\varepsilon< 4n-33< 4\varepsilon n+6\varepsilon
\\&\iff{1}
\frac{33-6\varepsilon}{4(\varepsilon+1)}< n<\frac{33+6\varepsilon}{4(1-\varepsilon)}
\end{align*}
כאשר במעבר האחרון חילקנו ב-$(1-\varepsilon)$
ולכן נניח כי הוא חיובי, כלומר כי
$\varepsilon<1$.
ניתן לראות כי אם נבחר
$\varepsilon=\frac12$
למשל, התנאי האחרון שקיבלנו מתקיים רק עבור מספר סופי של
$n$-ים
בסתירה לכך שהוא אמור להתקיים החל
מ-$N$
מסוים.
\end{solbox}

\begin{solbox}
דרך ב׳:
נראה שקיים
$\varepsilon>0$
כך שלכל
$N\in\mathbb{N}$
קיים
$n>N$
אשר יקיים
$$.\left|\frac{5n-12}{2n+3}-\frac{3}{2}\right|\geq\varepsilon$$
נפשט את הביטוי בתוך הערך המוחלט ונקבל
$$.\left| a_n-\frac{3}{2}\right|=\left|\frac{4n-33}{4n+6}\right|\geq\varepsilon$$
מסובך לעבוד עם ערכים מוחלטים, לכן ננסה "להיפטר" מהערך המוחלט, לשם כך נשים לב כי המכנה הוא חיובי ואילו המונה הוא חיובי רק כאשר
$n\geq 9$.
לכן נניח כי
$n\geq 9$
ונמשיך
$$.\left|\frac{5n-12}{2n+3}-\frac{3}{2}\right|=\frac{4n-33}{4n+6}=1-\frac{39}{4n+6}\geq 1-\frac{39}{4\cdot 9+6}=\frac{1}{14}$$
איך התוצאה שהתקבלה ניתנת לפירוש?
אם נבחר
$0<\varepsilon<\frac{1}{14}$
(לא משנה איזה)
אז הראנו כי \textbf{לכל}
$n\ge9$
מתקיים
$$\left|a_n-\frac{3}{2}\right|\geq\frac{1}{14}>\varepsilon$$
לכן לא קיים
$N$
עבורו לכל
$n>N$
מתקיים
$|a_n-\frac{3}{2}|<\varepsilon$.
זה מוכיח כי
$\frac{3}{2}$
אינו גבול של הסדרה.
\end{solbox}
\eMs
\end{enumerate}
